% --- Archon physics formalization source begin ---
% archon:physics
% archon:covers PhyXMiniProblems/problem_phyx_mini_0922.lean
% archon:source-report reports/phyx_mini/problem_phyx_mini_0922.source.json

\chapter{Physics problem phyx\_mini\_0922}
\label{ch:PhyXMiniProblems_problem_phyx_mini_0922}

\paragraph{Problem source.}
\#\# Physical scenario

Assume that \textbackslash{}( B \textbackslash{}) is uniform across a cross section (that is, neglect the variation of \textbackslash{}( B \textbackslash{}) with distance from the toroid axis).

\#\# Question

Determine the self-inductance of a toroidal solenoid with cross-sectional area \textbackslash{}( A \textbackslash{}) and mean radius \textbackslash{}( r \textbackslash{}), closely wound with \textbackslash{}( N \textbackslash{}) turns of wire on a nonmagnetic core.

\#\# Answer choices

- A: A: 20μH

- B: B: 80μH

- C: C: 40μH

- D: D: 45μH

\#\# Figure caption (auxiliary; use the image as primary evidence)

The image shows a diagram of a solenoid with a toroidal shape. Key elements include:

1. **Coil Turns**: Illustrated with a few loops around the toroid, labeled to indicate that the actual number of turns is \textbackslash{}( N \textbackslash{}), though only a few are shown for simplicity.

2. **Arrows**: Purple arrows on the coils indicate the direction of electric current (\textbackslash{}( i \textbackslash{})) flowing through the solenoid.

3. **Toroidal Core**: The core is depicted as a doughnut-shaped object, around which the coil is wound.

4. **Labels**:
   - Point \textbackslash{}( A \textbackslash{}) is marked on the toroid.
   - Radius \textbackslash{}( r \textbackslash{}) is indicated by a dashed arrow spanning from the center of the toroid to the inside of the coil.

5. **Text**: 
   - "Number of turns = \textbackslash{}( N \textbackslash{}) (only a few are shown)" explains the representation of the coil turns.
   - \textbackslash{}( i \textbackslash{}) denotes the electric current flowing into the solenoid.

This figure typically represents a toroidal inductor commonly used in physics and engineering to illustrate magnetic field concepts.

\#\# Dataset metadata

Electromagnetic Induction; Physical Model Grounding Reasoning, Multi-Formula Reasoning

\paragraph{Recorded answer/context.}
C: C: 40μH

\paragraph{Figure/image path.}
/root/proposal\_for\_physic/science-mango/phyx\_mini\_run/phyx\_data/test\_image/922.png

\paragraph{Formalization target.}
create a compiling Lean file with sorry bodies at `PhyXMiniProblems/problem\_phyx\_mini\_0922.lean`. The Lean declarations must preserve the physical quantities, dimensions or dimensional roles, figure labels, governing-law hypotheses, and final relation expressed by this problem.
Use Mathlib/Physlib names found through LeanExplore where available. If a physics API is missing, introduce faithful local abstractions rather than scalar placeholder aliases.

\begin{theorem}[Physics formalization target]
\label{thm:physics:phyx_mini_0922:target}
\lean{PhyXMiniProblems.ProblemPhyXMini0922.selfInductance_of_uniform_toroidal_solenoid}
\uses{def:physics:phyx-mini-0922:phyxminiproblems-problemphyxmini0922-lengthinmeters, def:physics:phyx-mini-0922:phyxminiproblems-problemphyxmini0922-areainsquaremeters, def:physics:phyx-mini-0922:phyxminiproblems-problemphyxmini0922-permeabilityinhenriespermeter, def:physics:phyx-mini-0922:phyxminiproblems-problemphyxmini0922-inductanceinhenries, def:physics:phyx-mini-0922:phyxminiproblems-problemphyxmini0922-toroidalsolenoidsetup, def:physics:phyx-mini-0922:phyxminiproblems-problemphyxmini0922-matchesprimaryfigure, def:physics:phyx-mini-0922:phyxminiproblems-problemphyxmini0922-matchesproblemscenario, def:physics:phyx-mini-0922:phyxminiproblems-problemphyxmini0922-hasphysicalparameters, def:physics:phyx-mini-0922:phyxminiproblems-problemphyxmini0922-satisfiestoroidalsolenoidlaws}
For a thin, closely wound toroidal solenoid with nonmagnetic core and magnetic
field uniform across its cross section,
\[
 L=\frac{\mu_0N^2A}{2\pi r}.
\]
\end{theorem}
\begin{proof}
Write \(B\) for the uniform field readout at the reference cross-section
point, \(i\) for the positive test current, and \(\mu_0,A,r,N\) for the
corresponding physical readouts.  The nonmagnetic-core hypothesis sets the
relative permeability to \(1\).  Ampère's law gives
\[
 B(2\pi r)=\mu_0Ni.
\]
Since \(r>0\) and \(\pi>0\), this determines
\(B=\mu_0Ni/(2\pi r)\).  Uniformity gives flux per turn
\(\Phi=BA\), and the linkage law gives
\(\lambda=N\Phi\).  The definition of self-inductance also gives
\(\lambda=Li\).  Substitute the first three identities into the last one and
cancel the positive current \(i\); the result is
\(L=\mu_0N^2A/(2\pi r)\).
\end{proof}

% --- Archon declaration topology begin ---
\section{Declaration topology}
\begin{definition}[electric Current Dimension]
  \label{def:physics:phyx-mini-0922:phyxminiproblems-problemphyxmini0922-electriccurrentdimension}
  \lean{PhyXMiniProblems.ProblemPhyXMini0922.electricCurrentDimension}
  Electric-current dimension Q T⁻¹
\end{definition}
\begin{proof}
  This is the declared model definition of electric Current Dimension.
\end{proof}
\begin{definition}[area Dimension]
  \label{def:physics:phyx-mini-0922:phyxminiproblems-problemphyxmini0922-areadimension}
  \lean{PhyXMiniProblems.ProblemPhyXMini0922.areaDimension}
  Cross-sectional-area dimension L²
\end{definition}
\begin{proof}
  This is the declared model definition of area Dimension.
\end{proof}
\begin{definition}[magnetic Flux Density Dimension]
  \label{def:physics:phyx-mini-0922:phyxminiproblems-problemphyxmini0922-magneticfluxdensitydimension}
  \lean{PhyXMiniProblems.ProblemPhyXMini0922.magneticFluxDensityDimension}
  Magnetic-flux-density dimension M T⁻¹ Q⁻¹ (tesla)
\end{definition}
\begin{proof}
  This is the declared model definition of magnetic Flux Density Dimension.
\end{proof}
\begin{definition}[magnetic Flux Dimension]
  \label{def:physics:phyx-mini-0922:phyxminiproblems-problemphyxmini0922-magneticfluxdimension}
  \lean{PhyXMiniProblems.ProblemPhyXMini0922.magneticFluxDimension}
  \uses{def:physics:phyx-mini-0922:phyxminiproblems-problemphyxmini0922-areadimension, def:physics:phyx-mini-0922:phyxminiproblems-problemphyxmini0922-magneticfluxdensitydimension}
  Magnetic-flux dimension M L² T⁻¹ Q⁻¹ (weber)
\end{definition}
\begin{proof}
  This is the declared model definition of magnetic Flux Dimension.
\end{proof}
\begin{definition}[magnetic Permeability Dimension]
  \label{def:physics:phyx-mini-0922:phyxminiproblems-problemphyxmini0922-magneticpermeabilitydimension}
  \lean{PhyXMiniProblems.ProblemPhyXMini0922.magneticPermeabilityDimension}
  Magnetic-permeability dimension M L Q⁻² (henry per metre)
\end{definition}
\begin{proof}
  This is the declared model definition of magnetic Permeability Dimension.
\end{proof}
\begin{definition}[inductance Dimension]
  \label{def:physics:phyx-mini-0922:phyxminiproblems-problemphyxmini0922-inductancedimension}
  \lean{PhyXMiniProblems.ProblemPhyXMini0922.inductanceDimension}
  Inductance dimension M L² Q⁻² (henry)
\end{definition}
\begin{proof}
  This is the declared model definition of inductance Dimension.
\end{proof}
\begin{definition}[Length Quantity]
  \label{def:physics:phyx-mini-0922:phyxminiproblems-problemphyxmini0922-lengthquantity}
  \lean{PhyXMiniProblems.ProblemPhyXMini0922.LengthQuantity}
  A nonnegative, unit-independent physical length
\end{definition}
\begin{proof}
  This is the declared model definition of Length Quantity.
\end{proof}
\begin{definition}[Area Quantity]
  \label{def:physics:phyx-mini-0922:phyxminiproblems-problemphyxmini0922-areaquantity}
  \lean{PhyXMiniProblems.ProblemPhyXMini0922.AreaQuantity}
  \uses{def:physics:phyx-mini-0922:phyxminiproblems-problemphyxmini0922-areadimension}
  A nonnegative, unit-independent physical area
\end{definition}
\begin{proof}
  This is the declared model definition of Area Quantity.
\end{proof}
\begin{definition}[Electric Current Quantity]
  \label{def:physics:phyx-mini-0922:phyxminiproblems-problemphyxmini0922-electriccurrentquantity}
  \lean{PhyXMiniProblems.ProblemPhyXMini0922.ElectricCurrentQuantity}
  \uses{def:physics:phyx-mini-0922:phyxminiproblems-problemphyxmini0922-electriccurrentdimension}
  A nonnegative winding-current magnitude
\end{definition}
\begin{proof}
  This is the declared model definition of Electric Current Quantity.
\end{proof}
\begin{definition}[Magnetic Flux Density Quantity]
  \label{def:physics:phyx-mini-0922:phyxminiproblems-problemphyxmini0922-magneticfluxdensityquantity}
  \lean{PhyXMiniProblems.ProblemPhyXMini0922.MagneticFluxDensityQuantity}
  \uses{def:physics:phyx-mini-0922:phyxminiproblems-problemphyxmini0922-magneticfluxdensitydimension}
  A nonnegative tangential magnetic-flux-density magnitude
\end{definition}
\begin{proof}
  This is the declared model definition of Magnetic Flux Density Quantity.
\end{proof}
\begin{definition}[Magnetic Flux Quantity]
  \label{def:physics:phyx-mini-0922:phyxminiproblems-problemphyxmini0922-magneticfluxquantity}
  \lean{PhyXMiniProblems.ProblemPhyXMini0922.MagneticFluxQuantity}
  \uses{def:physics:phyx-mini-0922:phyxminiproblems-problemphyxmini0922-magneticfluxdimension}
  A nonnegative magnetic flux or flux-linkage magnitude
\end{definition}
\begin{proof}
  This is the declared model definition of Magnetic Flux Quantity.
\end{proof}
\begin{definition}[Magnetic Permeability Quantity]
  \label{def:physics:phyx-mini-0922:phyxminiproblems-problemphyxmini0922-magneticpermeabilityquantity}
  \lean{PhyXMiniProblems.ProblemPhyXMini0922.MagneticPermeabilityQuantity}
  \uses{def:physics:phyx-mini-0922:phyxminiproblems-problemphyxmini0922-magneticpermeabilitydimension}
  A nonnegative magnetic permeability
\end{definition}
\begin{proof}
  This is the declared model definition of Magnetic Permeability Quantity.
\end{proof}
\begin{definition}[Inductance Quantity]
  \label{def:physics:phyx-mini-0922:phyxminiproblems-problemphyxmini0922-inductancequantity}
  \lean{PhyXMiniProblems.ProblemPhyXMini0922.InductanceQuantity}
  \uses{def:physics:phyx-mini-0922:phyxminiproblems-problemphyxmini0922-inductancedimension}
  A nonnegative self-inductance
\end{definition}
\begin{proof}
  This is the declared model definition of Inductance Quantity.
\end{proof}
\begin{definition}[nonnegative SIReadout]
  \label{def:physics:phyx-mini-0922:phyxminiproblems-problemphyxmini0922-nonnegativesireadout}
  \lean{PhyXMiniProblems.ProblemPhyXMini0922.nonnegativeSIReadout}
  Read any nonnegative dimensionful quantity in coherent SI units
\end{definition}
\begin{proof}
  This is the declared model definition of nonnegative SIReadout.
\end{proof}
\begin{definition}[length In Meters]
  \label{def:physics:phyx-mini-0922:phyxminiproblems-problemphyxmini0922-lengthinmeters}
  \lean{PhyXMiniProblems.ProblemPhyXMini0922.lengthInMeters}
  \uses{def:physics:phyx-mini-0922:phyxminiproblems-problemphyxmini0922-lengthquantity, def:physics:phyx-mini-0922:phyxminiproblems-problemphyxmini0922-nonnegativesireadout}
  Metre readout of a physical length
\end{definition}
\begin{proof}
  This is the declared model definition of length In Meters.
\end{proof}
\begin{definition}[area In Square Meters]
  \label{def:physics:phyx-mini-0922:phyxminiproblems-problemphyxmini0922-areainsquaremeters}
  \lean{PhyXMiniProblems.ProblemPhyXMini0922.areaInSquareMeters}
  \uses{def:physics:phyx-mini-0922:phyxminiproblems-problemphyxmini0922-areaquantity, def:physics:phyx-mini-0922:phyxminiproblems-problemphyxmini0922-nonnegativesireadout}
  Square-metre readout of a physical area
\end{definition}
\begin{proof}
  This is the declared model definition of area In Square Meters.
\end{proof}
\begin{definition}[current In Amperes]
  \label{def:physics:phyx-mini-0922:phyxminiproblems-problemphyxmini0922-currentinamperes}
  \lean{PhyXMiniProblems.ProblemPhyXMini0922.currentInAmperes}
  \uses{def:physics:phyx-mini-0922:phyxminiproblems-problemphyxmini0922-electriccurrentquantity, def:physics:phyx-mini-0922:phyxminiproblems-problemphyxmini0922-nonnegativesireadout}
  Ampere readout of a winding-current magnitude
\end{definition}
\begin{proof}
  This is the declared model definition of current In Amperes.
\end{proof}
\begin{definition}[magnetic Flux Density In Teslas]
  \label{def:physics:phyx-mini-0922:phyxminiproblems-problemphyxmini0922-magneticfluxdensityinteslas}
  \lean{PhyXMiniProblems.ProblemPhyXMini0922.magneticFluxDensityInTeslas}
  \uses{def:physics:phyx-mini-0922:phyxminiproblems-problemphyxmini0922-magneticfluxdensityquantity, def:physics:phyx-mini-0922:phyxminiproblems-problemphyxmini0922-nonnegativesireadout}
  Tesla readout of a tangential magnetic-flux-density magnitude
\end{definition}
\begin{proof}
  This is the declared model definition of magnetic Flux Density In Teslas.
\end{proof}
\begin{definition}[magnetic Flux In Webers]
  \label{def:physics:phyx-mini-0922:phyxminiproblems-problemphyxmini0922-magneticfluxinwebers}
  \lean{PhyXMiniProblems.ProblemPhyXMini0922.magneticFluxInWebers}
  \uses{def:physics:phyx-mini-0922:phyxminiproblems-problemphyxmini0922-magneticfluxquantity, def:physics:phyx-mini-0922:phyxminiproblems-problemphyxmini0922-nonnegativesireadout}
  Weber readout of magnetic flux or flux linkage
\end{definition}
\begin{proof}
  This is the declared model definition of magnetic Flux In Webers.
\end{proof}
\begin{definition}[permeability In Henries Per Meter]
  \label{def:physics:phyx-mini-0922:phyxminiproblems-problemphyxmini0922-permeabilityinhenriespermeter}
  \lean{PhyXMiniProblems.ProblemPhyXMini0922.permeabilityInHenriesPerMeter}
  \uses{def:physics:phyx-mini-0922:phyxminiproblems-problemphyxmini0922-magneticpermeabilityquantity, def:physics:phyx-mini-0922:phyxminiproblems-problemphyxmini0922-nonnegativesireadout}
  Henry-per-metre readout of magnetic permeability
\end{definition}
\begin{proof}
  This is the declared model definition of permeability In Henries Per Meter.
\end{proof}
\begin{definition}[inductance In Henries]
  \label{def:physics:phyx-mini-0922:phyxminiproblems-problemphyxmini0922-inductanceinhenries}
  \lean{PhyXMiniProblems.ProblemPhyXMini0922.inductanceInHenries}
  \uses{def:physics:phyx-mini-0922:phyxminiproblems-problemphyxmini0922-inductancequantity, def:physics:phyx-mini-0922:phyxminiproblems-problemphyxmini0922-nonnegativesireadout}
  Henry readout of self-inductance
\end{definition}
\begin{proof}
  This is the declared model definition of inductance In Henries.
\end{proof}
\begin{definition}[Figure Feature]
  \label{def:physics:phyx-mini-0922:phyxminiproblems-problemphyxmini0922-figurefeature}
  \lean{PhyXMiniProblems.ProblemPhyXMini0922.FigureFeature}
  Physical features explicitly visible in image 922.png
\end{definition}
\begin{proof}
  This is the declared model definition of Figure Feature.
\end{proof}
\begin{definition}[Figure Label]
  \label{def:physics:phyx-mini-0922:phyxminiproblems-problemphyxmini0922-figurelabel}
  \lean{PhyXMiniProblems.ProblemPhyXMini0922.FigureLabel}
  Mathematical labels printed in the primary figure
\end{definition}
\begin{proof}
  This is the declared model definition of Figure Label.
\end{proof}
\begin{definition}[Figure Label Role]
  \label{def:physics:phyx-mini-0922:phyxminiproblems-problemphyxmini0922-figurelabelrole}
  \lean{PhyXMiniProblems.ProblemPhyXMini0922.FigureLabelRole}
  Physical roles of the labels printed in the primary figure
\end{definition}
\begin{proof}
  This is the declared model definition of Figure Label Role.
\end{proof}
\begin{definition}[Toroidal Solenoid Figure]
  \label{def:physics:phyx-mini-0922:phyxminiproblems-problemphyxmini0922-toroidalsolenoidfigure}
  \lean{PhyXMiniProblems.ProblemPhyXMini0922.ToroidalSolenoidFigure}
  \uses{def:physics:phyx-mini-0922:phyxminiproblems-problemphyxmini0922-figurefeature, def:physics:phyx-mini-0922:phyxminiproblems-problemphyxmini0922-figurelabel, def:physics:phyx-mini-0922:phyxminiproblems-problemphyxmini0922-figurelabelrole}
  Literal and qualitative evidence transcribed from image 922.png
\end{definition}
\begin{proof}
  This is the declared model definition of Toroidal Solenoid Figure.
\end{proof}
\begin{definition}[Core Material]
  \label{def:physics:phyx-mini-0922:phyxminiproblems-problemphyxmini0922-corematerial}
  \lean{PhyXMiniProblems.ProblemPhyXMini0922.CoreMaterial}
  Magnetic behavior of the toroidal core material
\end{definition}
\begin{proof}
  This is the declared model definition of Core Material.
\end{proof}
\begin{definition}[Winding Style]
  \label{def:physics:phyx-mini-0922:phyxminiproblems-problemphyxmini0922-windingstyle}
  \lean{PhyXMiniProblems.ProblemPhyXMini0922.WindingStyle}
  Whether adjacent turns realize the closely wound idealization
\end{definition}
\begin{proof}
  This is the declared model definition of Winding Style.
\end{proof}
\begin{definition}[Current Orientation]
  \label{def:physics:phyx-mini-0922:phyxminiproblems-problemphyxmini0922-currentorientation}
  \lean{PhyXMiniProblems.ProblemPhyXMini0922.CurrentOrientation}
  Sign convention for current relative to the purple arrows in the figure
\end{definition}
\begin{proof}
  This is the declared model definition of Current Orientation.
\end{proof}
\begin{definition}[Toroidal Solenoid Setup]
  \label{def:physics:phyx-mini-0922:phyxminiproblems-problemphyxmini0922-toroidalsolenoidsetup}
  \lean{PhyXMiniProblems.ProblemPhyXMini0922.ToroidalSolenoidSetup}
  \uses{def:physics:phyx-mini-0922:phyxminiproblems-problemphyxmini0922-lengthquantity, def:physics:phyx-mini-0922:phyxminiproblems-problemphyxmini0922-areaquantity, def:physics:phyx-mini-0922:phyxminiproblems-problemphyxmini0922-electriccurrentquantity, def:physics:phyx-mini-0922:phyxminiproblems-problemphyxmini0922-magneticfluxdensityquantity, def:physics:phyx-mini-0922:phyxminiproblems-problemphyxmini0922-magneticfluxquantity, def:physics:phyx-mini-0922:phyxminiproblems-problemphyxmini0922-magneticpermeabilityquantity, def:physics:phyx-mini-0922:phyxminiproblems-problemphyxmini0922-inductancequantity, def:physics:phyx-mini-0922:phyxminiproblems-problemphyxmini0922-toroidalsolenoidfigure, def:physics:phyx-mini-0922:phyxminiproblems-problemphyxmini0922-corematerial, def:physics:phyx-mini-0922:phyxminiproblems-problemphyxmini0922-windingstyle, def:physics:phyx-mini-0922:phyxminiproblems-problemphyxmini0922-currentorientation}
  Independent physical quantities of the toroidal solenoid. The abstract CrossSectionPoint type lets field uniformity be stated without replacing the cross section by a single scalar sample. In particular, selfInductance is not defined from either the requested formula or an answer choice
\end{definition}
\begin{proof}
  This is the declared model definition of Toroidal Solenoid Setup.
\end{proof}
\begin{definition}[Matches Primary Figure]
  \label{def:physics:phyx-mini-0922:phyxminiproblems-problemphyxmini0922-matchesprimaryfigure}
  \lean{PhyXMiniProblems.ProblemPhyXMini0922.MatchesPrimaryFigure}
  \uses{def:physics:phyx-mini-0922:phyxminiproblems-problemphyxmini0922-figurefeature, def:physics:phyx-mini-0922:phyxminiproblems-problemphyxmini0922-figurelabel, def:physics:phyx-mini-0922:phyxminiproblems-problemphyxmini0922-toroidalsolenoidsetup}
  All labels and qualitative features read from the primary bitmap
\end{definition}
\begin{proof}
  This is the declared model definition of Matches Primary Figure.
\end{proof}
\begin{definition}[Matches Problem Scenario]
  \label{def:physics:phyx-mini-0922:phyxminiproblems-problemphyxmini0922-matchesproblemscenario}
  \lean{PhyXMiniProblems.ProblemPhyXMini0922.MatchesProblemScenario}
  \uses{def:physics:phyx-mini-0922:phyxminiproblems-problemphyxmini0922-toroidalsolenoidsetup}
  Qualitative conditions stated in the problem text
\end{definition}
\begin{proof}
  This is the declared model definition of Matches Problem Scenario.
\end{proof}
\begin{definition}[Has Physical Parameters]
  \label{def:physics:phyx-mini-0922:phyxminiproblems-problemphyxmini0922-hasphysicalparameters}
  \lean{PhyXMiniProblems.ProblemPhyXMini0922.HasPhysicalParameters}
  \uses{def:physics:phyx-mini-0922:phyxminiproblems-problemphyxmini0922-lengthinmeters, def:physics:phyx-mini-0922:phyxminiproblems-problemphyxmini0922-areainsquaremeters, def:physics:phyx-mini-0922:phyxminiproblems-problemphyxmini0922-currentinamperes, def:physics:phyx-mini-0922:phyxminiproblems-problemphyxmini0922-permeabilityinhenriespermeter, def:physics:phyx-mini-0922:phyxminiproblems-problemphyxmini0922-toroidalsolenoidsetup}
  Physical nondegeneracy assumptions. The positive test current is used only to identify the inductance through lambda = L i; it does not prescribe the requested value of L
\end{definition}
\begin{proof}
  This is the declared model definition of Has Physical Parameters.
\end{proof}
\begin{definition}[Satisfies Toroidal Solenoid Laws]
  \label{def:physics:phyx-mini-0922:phyxminiproblems-problemphyxmini0922-satisfiestoroidalsolenoidlaws}
  \lean{PhyXMiniProblems.ProblemPhyXMini0922.SatisfiesToroidalSolenoidLaws}
  \uses{def:physics:phyx-mini-0922:phyxminiproblems-problemphyxmini0922-lengthinmeters, def:physics:phyx-mini-0922:phyxminiproblems-problemphyxmini0922-areainsquaremeters, def:physics:phyx-mini-0922:phyxminiproblems-problemphyxmini0922-currentinamperes, def:physics:phyx-mini-0922:phyxminiproblems-problemphyxmini0922-magneticfluxdensityinteslas, def:physics:phyx-mini-0922:phyxminiproblems-problemphyxmini0922-magneticfluxinwebers, def:physics:phyx-mini-0922:phyxminiproblems-problemphyxmini0922-permeabilityinhenriespermeter, def:physics:phyx-mini-0922:phyxminiproblems-problemphyxmini0922-inductanceinhenries, def:physics:phyx-mini-0922:phyxminiproblems-problemphyxmini0922-toroidalsolenoidsetup}
  The modeling relations used for a thin, closely wound toroid. These are governing laws and intermediate physical relations; none states the requested closed form for self-inductance
\end{definition}
\begin{proof}
  This is the declared model definition of Satisfies Toroidal Solenoid Laws.
\end{proof}
\begin{definition}[Answer Choice]
  \label{def:physics:phyx-mini-0922:phyxminiproblems-problemphyxmini0922-answerchoice}
  \lean{PhyXMiniProblems.ProblemPhyXMini0922.AnswerChoice}
  Labels of the four printed multiple-choice answers
\end{definition}
\begin{proof}
  This is the declared model definition of Answer Choice.
\end{proof}
\begin{definition}[answer In Microhenries]
  \label{def:physics:phyx-mini-0922:phyxminiproblems-problemphyxmini0922-answerinmicrohenries}
  \lean{PhyXMiniProblems.ProblemPhyXMini0922.answerInMicrohenries}
  \uses{def:physics:phyx-mini-0922:phyxminiproblems-problemphyxmini0922-answerchoice}
  Printed answer magnitude in microhenries
\end{definition}
\begin{proof}
  This is the declared model definition of answer In Microhenries.
\end{proof}
\begin{definition}[answer In Henries]
  \label{def:physics:phyx-mini-0922:phyxminiproblems-problemphyxmini0922-answerinhenries}
  \lean{PhyXMiniProblems.ProblemPhyXMini0922.answerInHenries}
  \uses{def:physics:phyx-mini-0922:phyxminiproblems-problemphyxmini0922-answerchoice, def:physics:phyx-mini-0922:phyxminiproblems-problemphyxmini0922-answerinmicrohenries}
  Convert a printed microhenry magnitude to henries
\end{definition}
\begin{proof}
  This is the declared model definition of answer In Henries.
\end{proof}
\begin{definition}[recorded Answer Choice]
  \label{def:physics:phyx-mini-0922:phyxminiproblems-problemphyxmini0922-recordedanswerchoice}
  \lean{PhyXMiniProblems.ProblemPhyXMini0922.recordedAnswerChoice}
  \uses{def:physics:phyx-mini-0922:phyxminiproblems-problemphyxmini0922-answerchoice}
  Dataset metadata records answer choice C; it is not a governing law
\end{definition}
\begin{proof}
  This is the declared model definition of recorded Answer Choice.
\end{proof}
% --- Archon declaration topology end ---
% --- Archon physics formalization source end ---
